\section{Conclusions}
\label{conclusions}
Cities, with their high population density, account for most of the oil and other fats consumed for culinary purposes, whether for domestic use or services. In the current context of the fight against climate change and the achievement of the SDGs, cities play a critical role, as they are responsible for the proper treatment of waste and for trying to implement the necessary measures to combat global warming. In this sense, UCO can be a valuable ally, given its ease of reuse as biodiesel, which produces clean and sustainable energy. However, planning urban infrastructure for its collection using containers equipped for this purpose is not always easy. These types of facilities have often grown organically alongside cities, resulting in deficient access in terms of coverage and equity, which discourages greater collection.

In this paper, we introduced a multi-objective facility location formulation and a tailored evolutionary algorithm for the placement of UCO collection containers. The model integrates four planning criteria: the number of containers, the population left without capacity in Voronoi-based service areas, the spatial equity of that shortage, and the distance to the nearest container. We showed that Moran's index, a natural candidate for the equity criterion, is ill-posed as an objective: it is fixed by the number of observations in small solutions, nearly constant in realistic ones, and blind to the magnitude of the shortage. We replaced it with a measure on a fixed spatial support, the population-weighted dispersion of the unserved share between neighbourhoods, which is comparable across solutions of any size and grows only when the shortage is both unequal and spatially concentrated. We showed that this measure is robust to its parameters and ranks solutions differently from the classical non-spatial equity measures. Methodologically, we adapted NSGA-II to sets of locations with a discrete set of network sizes and three problem-specific operators: a partition crossover between parents of the same size with load-guided repair, a mutation that moves containers towards the people they serve, and an initialization with equal-population cells. An extensive study on real data from five Spanish cities showed that the size management and the load-guided resizing are essential to the quality of the fronts; that NSGA-III brings no improvement over NSGA-II, even with four objectives; that, in the same computing time, the proposed configuration improves the hypervolume of a mutation-driven one by 4--10\% and that of seven crossovers of the literature, in their original form, by 3--11\%; that same-size mating and load-guided repair are what make recombination work, for our crossover and for the others, so that crossovers that exchange large regions reach the level of ours once they have them; that the exact decomposition of two layouts into independent components, on which partition crossover rests in other problems, leaves one component with most of the differences, and imposing tiles works better; that the equal-population seeds alone are fairer and more accessible than today's layouts; and that adding distance as an objective costs little in the other objectives while covering the trade-off between equity and accessibility that the three-objective formulation leaves to chance. The most accessible fair layouts are within 6--9\% of the best known solutions of the corresponding $p$-median problem, whose optimum we bracketed within 6\% with a Lagrangian bound.

From an applied perspective, the fronts are directly usable by planners. With today's number of containers, every front contained layouts that are better than the current deployment in unserved population, spatial inequity (by a factor of 2 to 7) and mean distance at once, and layouts nearly as accessible as those of a single-objective optimizer but fairer. The fronts also make explicit what each additional container buys, and show that in cities with a strong shortage of capacity, such as Málaga and Sevilla, relocation alone cannot serve more people, but it can make the shortage fair.

We recognize the limitations that shape future research. First, no real cost data were available, so the size of the network stands for its cost; per-site installation and maintenance costs can replace it directly. Second, with four objectives most of the population is non-dominated, which weakens the pressure towards the extremes of the front: the fairest layouts are somewhat less fair than with three objectives. Larger populations, or archives per network size, could recover them at a higher computational cost. The proposed configuration improves most of the front but not its fairest end, which needs more generations or the comparison of inequity in each tile; an operator that handles the non-additive objective directly is an open question. Once mating and repair are in place, the differences between the best crossovers are within 1\%, so the remaining room lies less in the shape of the exchanged regions than in the mutation, the initialization and the treatment of inequity; partitions that follow the structure of the city, such as districts or the street network, are still worth testing. Third, distances are straight-line; walking distances on the street network would be more realistic. Fourth, the model assumes static demand and that citizens use the nearest container; data on container filling and on actual use would make it more realistic. Finally, combining location with collection routing would capture the operational cost of the network.

Taken together, our findings show that integrating spatial equity into multi-objective evolutionary optimization provides a flexible framework for siting urban waste collection infrastructure, provided the equity measure is well-posed for optimization. The methodological lessons presented here, equity measured on a fixed spatial support, problem-specific mutation and size management, and recombination that uses the locality of the objectives rather than blind recombination, and bounds on the classical single-objective problem to assess solution quality, carry over to other urban services that face similar trade-offs between cost, coverage and fairness. The method is deployed as an open-source service of a data space, so municipalities that publish their container inventories can use it directly.
